\documentclass[varwidth=16cm, border=2mm]{standalone}
\usepackage{amsmath}\usepackage{amssymb}\usepackage{ctex}
\usepackage{tasks}\usepackage{xcolor}\usepackage{graphicx}
\settasks{label=\Alph*., label-width=1.5em, item-indent=2em}
\begin{document}
\textbf{[0.6]} (2024·全国·高三对口高考) 已知椭圆 $\frac{x^2}{9} + y^2 = 1$，过左焦点 $F$ 作倾斜角为 $\frac{\pi}{6}$ 的直线交椭圆于 $A、B$ 两点，则弦 $AB$ 的长为 _____.

\vspace{0.2cm}\par\noindent\textcolor{red}{\textbf{【答案】} 2}
\par\noindent\textcolor{blue}{\textbf{【解析】} 1. **确定椭圆参数**
椭圆方程为 $\frac{x^2}{9} + y^2 = 1$。
对照标准方程 $\frac{x^2}{a^2} + \frac{y^2}{b^2} = 1$，可得 $a^2 = 9$， $b^2 = 1$。
所以 $a=3$， $b=1$。

2. **计算焦距**
$c^2 = a^2 - b^2 = 9 - 1 = 8$。
所以 $c = \sqrt{8} = 2\sqrt{2}$。
左焦点 $F$ 的坐标为 $(-c, 0) = (-2\sqrt{2}, 0)$。

3. **确定直线的方程**
直线过左焦点 $F(-2\sqrt{2}, 0)$，倾斜角为 $\frac{\pi}{6}$。
直线的斜率 $k = \tan(\frac{\pi}{6}) = \frac{\sqrt{3}}{3}$。
根据点斜式，直线的方程为 $y - 0 = \frac{\sqrt{3}}{3}(x - (-2\sqrt{2}))$。
即 $y = \frac{\sqrt{3}}{3}(x + 2\sqrt{2})$。

4. **联立直线与椭圆方程**
将直线方程代入椭圆方程 $\frac{x^2}{9} + y^2 = 1$:
$\frac{x^2}{9} + \left(\frac{\sqrt{3}}{3}(x + 2\sqrt{2})\right)^2 = 1$
$\frac{x^2}{9} + \frac{3}{9}(x + 2\sqrt{2})^2 = 1$
$\frac{x^2}{9} + \frac{1}{3}(x^2 + 4\sqrt{2}x + 8) = 1$
为了消去分母，方程两边同乘以 9:
$x^2 + 3(x^2 + 4\sqrt{2}x + 8) = 9$
$x^2 + 3x^2 + 12\sqrt{2}x + 24 = 9$
$4x^2 + 12\sqrt{2}x + 15 = 0$

5. **计算弦长**
设直线与椭圆交于 $A(x_1, y_1)$ 和 $B(x_2, y_2)$ 两点。
根据韦达定理，对于二次方程 $4x^2 + 12\sqrt{2}x + 15 = 0$，有：
$x_1 + x_2 = -\frac{12\sqrt{2}}{4} = -3\sqrt{2}$
$x_1 x_2 = \frac{15}{4}$
弦长公式为 $|AB| = \sqrt{(x_1 - x_2)^2 + (y_1 - y_2)^2}$。
由于 $y_1 = kx_1 + b'$ 且 $y_2 = kx_2 + b'$ (其中 $b' = \frac{2\sqrt{6}}{3}$)，所以 $y_1 - y_2 = k(x_1 - x_2)$。
将此关系代入弦长公式：
$|AB| = \sqrt{(x_1 - x_2)^2 + k^2(x_1 - x_2)^2} = \sqrt{(1+k^2)(x_1 - x_2)^2} = |x_1 - x_2|\sqrt{1+k^2}$。
首先计算 $(x_1 - x_2)^2$: 
$(x_1 - x_2)^2 = (x_1 + x_2)^2 - 4x_1 x_2 = (-3\sqrt{2})^2 - 4\left(\frac{15}{4}\right)$
$= (9 \times 2) - 15 = 18 - 15 = 3$。
所以 $|x_1 - x_2| = \sqrt{3}$。
又因为 $k = \frac{\sqrt{3}}{3}$，所以 $k^2 = \left(\frac{\sqrt{3}}{3}\right)^2 = \frac{3}{9} = \frac{1}{3}$。
代入弦长公式：
$|AB| = \sqrt{3} \cdot \sqrt{1 + \frac{1}{3}} = \sqrt{3} \cdot \sqrt{\frac{4}{3}} = \sqrt{3} \cdot \frac{2}{\sqrt{3}} = 2$。

**方法二：使用焦点弦长公式**
对于过焦点 $(c,0)$ 的弦，其倾斜角为 $\theta$，弦长公式为 $L = \frac{2ab^2}{a^2 - c^2 \cos^2\theta}$。
本题中 $a=3, b=1, c=2\sqrt{2}$，倾斜角 $\theta = \frac{\pi}{6}$。
$L = \frac{2 \times 3 \times 1^2}{3^2 - (2\sqrt{2})^2 \cos^2(\frac{\pi}{6})}$
$L = \frac{6}{9 - 8 \times \left(\frac{\sqrt{3}}{2}\right)^2}$
$L = \frac{6}{9 - 8 \times \frac{3}{4}}$
$L = \frac{6}{9 - 6}$
$L = \frac{6}{3} = 2$。
两种方法结果一致。

最终答案为 $2$。}


\end{document}
